\documentclass[11pt,letterpaper]{amsart}

\usepackage{amsmath}
\usepackage{graphicx}
\usepackage{amsmath,amssymb,amsfonts,amsthm,graphicx,amsopn,mathrsfs}
\usepackage[shortlabels]{enumitem}
\usepackage{tikz-cd}
\usepackage{tikz}
\usepackage{stmaryrd,url}
\usepackage{mathdots, verbatim,mathtools}
\usepackage{array,float}
\usepackage{fullpage}
\usepackage[utf8]{inputenc}
\usepackage{multirow}
\usepackage[normalem]{ulem}
\usepackage{hyperref}
\usepackage{cleveref}
\usepackage{algorithmic,algorithm}
\usepackage{caption}


\tikzstyle{nodal}=[circle,draw,fill=black,inner sep=0pt, minimum width=4pt]
\tikzstyle{half-fiber}=[rectangle,draw=black,thick,inner sep=0pt, minimum width=5pt, minimum height=5pt]
\tikzset{double distance = 2pt}

\newcommand{\Si}[1]{\textcolor{blue}{(Si: #1)}}
\newcommand{\Gi}[1]{\textcolor{red}{(Gi: #1)}}
\newcommand{\vect}[1]{\boldsymbol{#1}}


\DeclareMathOperator{\Hom}{Hom}
\DeclareMathOperator{\GL}{GL}
\DeclareMathOperator{\Ker}{Ker}
\DeclareMathOperator{\coker}{coker}
\DeclareMathOperator{\rk}{rk}
\DeclareMathOperator{\Aut}{Aut}
\DeclareMathOperator{\id}{id}
\DeclareMathOperator{\ord}{ord}
\DeclareMathOperator{\Img}{Im}
\DeclareMathOperator{\vol}{vol}
\DeclareMathOperator{\disc}{disc}
\DeclareMathOperator{\Div}{Div}
\DeclareMathOperator{\Pic}{Pic}
\DeclareMathOperator{\NS}{NS}
\DeclareMathOperator{\MW}{MW}
\DeclareMathOperator{\Or}{O}
\DeclareMathOperator{\Nef}{Nef}


\usepackage[sorting=nyc,doi=false,isbn=false,url=false,giveninits=true,maxbibnames=99]{biblatex}

\DeclareSortingTemplate{nyc}{
  \sort{
    \field{presort}
  }
  \sort[final]{
    \field{sortkey}
  }
  \sort{
    \field{sortname}
    \field{author}
    \field{editor}
    \field{translator}
    \field{sorttitle}
    \field{title}
  }
  \sort{
    \field{sortyear}
    \field{year}
  }
  \sort{\citeorder}
}

\addbibresource{mybib.bib}

\AtEveryBibitem{%
\ifentrytype{article}{
    \clearfield{number}
}{}
}
\renewbibmacro{in:}{%
  \ifentrytype{article}{}{\printtext{\bibstring{in}\intitlepunct}}}
\AtEveryBibitem{%
\ifentrytype{book}{
    \clearfield{pages}
}{}
}

\title{Virtually abelian symmetry groups of hyperbolic lattices}
\author{Simon Brandhorst}
\address{Simon Brandhorst:
Fakult\"at f\"ur Mathematik und Informatik, Universit\"at des Saarlandes, Campus E2.4, 66123 Saarbr\"ucken, Germany}
\email{brandhorst@math.uni-sb.de}

\author{Giacomo Mezzedimi}
\address{Giacomo Mezzedimi:
Mathematisches Institut, Universit\"at Bonn, Endenicher Allee 60, 53115 Bonn, Germany.}
\email{mezzedim@math.uni-bonn.de}


\thanks{S.B. is funded by the Deutsche Forschungsgemeinschaft (DFG, German Research Foundation) – Project-ID 286237555 – TRR 195.
Gefördert durch die Deutsche Forschungsgemeinschaft (DFG) – Projektnummer 286237555 – TRR 195.}


\date{\today}


\newcommand{\QQ}{\mathbb{Q}}
\newcommand{\RR}{\mathbb{R}}
\newcommand{\ZZ}{\mathbb{Z}}
\newcommand{\FF}{\mathbb{F}}
\newcommand{\CC}{\mathbb{C}}
\newcommand{\NN}{\mathbb{N}}
\newcommand{\PP}{\mathbb{P}}
\newcommand{\HH}{\mathbb{H}}
\newcommand{\Stabe}{\Aut(\mathcal{D}_L,\vect{e})}
\newcommand{\Stabee}{\Aut(\mathcal{D}_L,\vect{e'})}

\newcommand{\D}{\mathcal{D}}
\newcommand{\F}{\mathcal{F}}

\newtheorem{theorem}{Theorem}[section]
\newtheorem{proposition}[theorem]{Proposition}
\newtheorem{lemma}[theorem]{Lemma}
\newtheorem{corollary}[theorem]{Corollary}
\newtheorem{claim}[theorem]{Claim}
\theoremstyle{definition}
\newtheorem*{acknowledgments}{\textbf{Acknowledgments}}
\newtheorem*{conv}{\textbf{Conventions}}
\newtheorem{definition}[theorem]{Definition}
\theoremstyle{remark}
\newtheorem{remark}[theorem]{Remark}
\newtheorem{question}[theorem]{Question}
\newtheorem{example}[theorem]{Example}

\makeatletter
\let\@wraptoccontribs\wraptoccontribs
\makeatother

\contrib[with an appendix by]{Simon Brandhorst and Markus Kirschmer}
\address{Markus Kirschmer:
Fakult\"at f\"ur Mathematik, Universit\"at Bielefeld, Universit\"atsstra\"sse 25, 33501 Bielefeld, Germany.}
\email{mkirschm@math.uni-bielefeld.de}


\begin{document}

\maketitle
\begin{abstract}
    We give a classification of integral lattices with virtually abelian symmetry group. As a consequence, we complete the classification of K3 surfaces with virtually abelian automorphism group.
    In the appendix we formulate an algorithm for weak approximation in orthogonal groups and use it to determine if two indefinite lattices are isometric. 
\end{abstract}


\section{Introduction}

According to the Tits alternative \cite{tits}, any finitely generated subgroup of isometries of a hyperbolic lattice $L$ is either virtually solvable or it contains a free, non-abelian subgroup, where we say that a property of a group holds \emph{virtually} if it holds for a finite index subgroup. In view of our geometric applications, we focus on a specific subgroup of isometries of $L$, namely its \emph{symmetry group}. Let $\HH_L$ be the hyperbolic space associated to the hyperbolic lattice $L$. The \emph{Weyl group} $W(L)\subseteq \Or^+(L)$ is the subgroup generated by reflections into $(-2)$-roots of $L$, and we denote by $\mathcal{D}_L$ the closure of a fundamental domain for the action of $W(L)$ on the positive cone of $L$. Then the symmetry group of $L$ can be equivalently defined as the quotient $\Aut(\mathcal{D}_L) \coloneqq \Or^+(L)/W(L)$, or as the subgroup of $\Or^+(L)$ preserving $\mathcal{D}_L$.

Basic results in hyperbolic geometry show that the symmetry group $\Aut(\mathcal{D}_L)$ is virtually solvable if and only if it is virtually abelian, and this is the case exactly when $\Aut(\mathcal{D}_L)$ is \emph{elementary}, i.e. it has a finite orbit in $\HH_L$ (cf. \cite[Theorems~5.5.9~and~5.5.10]{ratcliffe}). There are three types of elementary groups of isometries in $\Or^+(L)$ (cf. \cite[§5.5]{ratcliffe}):
\begin{enumerate}
\item Elementary groups \emph{of elliptic type} have a fixed point in the interior of $\HH_L$, and they are finite;
\item Elementary groups \emph{of parabolic type} have precisely one fixed point in the boundary $\partial \HH_L$, and they have zero entropy (cf. \cite[Theorem~3.7]{brandhorst.mezzedimi});
\item Elementary groups \emph{of hyperbolic type} have two fixed points in $\partial \HH_L$, and they are virtually cyclic.
\end{enumerate}

The main goal of this paper is to complete the classification of virtually solvable (or abelian) symmetry groups of hyperbolic lattices. In rank $\le 2$, it is well known that the symmetry group $\Aut(\mathcal{D}_L)$ is either finite or elementary of hyperbolic type (see e.g. \cite{galluzzi.lombardo.peters}). 
Hyperbolic lattices of rank $\ge 3$ with finite symmetry group were classified by Nikulin \cite{nikulin.finite.aut.greater.5, nikulin.finite.aut.3} and Vinberg \cite{vinberg.finite.aut.4}, while lattices with elementary symmetry groups of parabolic type were recently classified in \cite{brandhorst.mezzedimi} (and independently by X. Yu in \cite{yu.entropy} in the case of Picard lattices of K3 surfaces).
By Nikulin \cite[Theorem 1]{nikulin.interesting}, the number of isomorphism classes of lattices of a fixed rank $n\geq 3$ with elementary symmetry group of hyperbolic type is finite. 
%Therefore, it should be possible to classify them.
Our main result is the following:

\begin{theorem} \label{thm:main}
A hyperbolic lattice of rank $\ge 3$ has a virtually solvable symmetry group if and only if it belongs to an explicit list of $360$ lattices. More precisely, $\Aut(\mathcal{D}_L)$ is:
\begin{enumerate}
    \item finite if and only if $L$ belongs to Nikulin--Vinberg's list of $118$ lattices. In particular $\rk(L)\le 19$.
    \item elementary of parabolic type if and only if $L$ belongs to the list in \cite{brandhorst.mezzedimi} of $194$ Borcherds lattices. In particular $\rk(L)\le 26$.
    \item elementary of hyperbolic type if and only if $L$ belongs to the list of $48$ lattices in \Cref{thm:classification.3} and \Cref{thm:classification.4}. In particular $\rk(L)\le 4$.
\end{enumerate}
\end{theorem}

Thanks to the Torelli theorem, the automorphism group of a K3 surface $X$ over an algebraically closed field $k$ can be identified, up to a finite group, with the symmetry group of its Picard lattice $\Pic(X)$. Therefore a byproduct of the classification in \Cref{thm:main} is the following:

\begin{corollary}
Let $X$ be a K3 surface over an algebraically closed field $k$. The automorphism group $\Aut(X)$ is virtually solvable if and only if either $\rho(X)=\rk(\Pic(X))\le 2$, or $\Pic(X)$ belongs to an explicit list of $359$ lattices of rank $\ge 3$. If this is the case, the rank of a maximal abelian subgroup of $\Aut(X)$ is at most $8$.
\end{corollary}

K3 surfaces with an elementary automorphism group of elliptic or parabolic type are the so-called K3 surfaces \emph{of zero entropy}, and they were studied in our previous paper \cite{brandhorst.mezzedimi}. From a dynamical point of view, K3 surfaces with an elementary automorphism group of hyperbolic type are more interesting: all its automorphisms of infinite order have positive entropy \cite[Theorem~5.5.8]{ratcliffe}, and therefore its elliptic fibrations (if they exist at all) have finite Mordell--Weil group (see \Cref{lem:detect_virtually_abelian}).

\subsection*{Strategy and outline}
It follows from a result of Nikulin \cite[Theorem 9.1.1]{nikulin.factor.groups} that, if a hyperbolic lattice of rank $\ge 5$ has an elementary symmetry group, then it is either of elliptic or parabolic type (cf. \cite[Theorem~3.9]{brandhorst.mezzedimi}). Therefore, the core of our classification is in rank $3$ and $4$. Let $L$ by a hyperbolic lattice of rank $3$ or $4$ with an elementary symmetry group of hyperbolic type. Then it is $2$-reflective of hyperbolic type in the sense of \cite[Definition 2]{nikulin.interesting}. By \cite[Lemma~1]{nikulin.interesting} we know that $L$ admits a $\QQ$-basis of $(-2)$-roots with bounded intersection multiplicity. This makes the classification in \Cref{thm:main} into a finite, though impractically long, task of checking all the possible lattices. Therefore, we develop three tests, the \emph{overlattice test} (see \Cref{lem:overlattice}), the \emph{sublattice test} (see \Cref{alg:sublattice}) and the \emph{cusp test} (see \Cref{alg:cusp test}) that allow us to exclude all but few lattices. We then employ the first author's implementation in OSCAR \cite{OSCAR} of Borcherds' method \cite{shimada.borcherds} as formulated by Shimada to obtain the classification in \Cref{thm:main}. Our computations also confirm the correctness of Nikulin's and Vinberg's classifications of hyperbolic lattices with finite symmetry group and rank $3$ and $4$ in \cite{nikulin.finite.aut.3} and \cite{vinberg.finite.aut.4}. 

Let us briefly outline the contents of the paper. In \Cref{sec:preliminaries} we recall some known properties of hyperbolic lattices with an elementary symmetry group of hyperbolic type, and we prove the simple, but crucial, \Cref{prop:easy_fundamental_chamber}. In \Cref{sec:tests} we introduce the three tests discussed above, and then in \Cref{sec:classification} we apply them to obtain the final classification. 
In \Cref{appendix} we give an algorithm realising the weak approximation theorem for orthogonal groups and use it to determine if two indefinite lattices of rank at least $3$ are isometric. The algorithm is used in our classification to reduce computation times and to ensure that the final result consists of pairwise non-isometric lattices. 

\subsection*{Acknowledgments}
We warmly thank Serge Cantat for helpful discussions and comments.  This work started during the conference ``Birational geometry and dynamics'' at the SwissMAP Research Station Les Diablerets. We thank the organizers of the conference and the staff of the station for their hospitality.


\section{Preliminaries} \label{sec:preliminaries}

We keep notations as in our previous article \cite{brandhorst.mezzedimi}. We refer to \cite[§2]{brandhorst.mezzedimi} for the basic definitions about lattices, and only recall the notions that will be strictly necessary for our purposes.

Let $L$ be a hyperbolic lattice, and denote by $\HH_L$ the hyperbolic space associated to $L$. We choose a (closed) fundamental chamber $\D_L$ for the action of the Weyl group $W(L)$ on $L_\RR:=L \otimes \RR$, and, by a slight abuse of notation, we denote again by $\D_L$ the image of $\D_L$ in $\HH_L$. The chamber $\D_L$ is a locally polyhedral convex cone, whose walls are given by $r^\perp$, where $r$ runs along the simple $(-2)$-roots of $L$. The chambers $\{w(\D_L)\}_{w\in W(L)}$ are called \emph{Weyl chambers} of $L$, and they form a tessellation of $\HH_L$.
If $L'$ is another lattice in $L\otimes \QQ$ of full rank, then we always choose $\D_{L'}$ such that $\D_{L'}^\circ\cap \D_{L}^\circ\ne\emptyset$ (note that if $L'$ is of smaller rank, this may not be possible).

The group $\Aut(\D_L)$ of isometries of $L$ preserving the fundamental chamber $\D_L$ is called the \emph{symmetry group} of $L$. By construction it coincides with the quotient $\Or^+(L)/W(L)$. The symmetry group $\Aut(\D_L)$ is \emph{elementary} if it has a finite orbit in $\HH_L$. If $\Aut(\D_L)$ is infinite, we say that it is \emph{elementary of hyperbolic type} if it has two fixed points in the boundary $\partial \HH_L$ (cf. \cite[§5.5]{ratcliffe}). In this case $\Aut(\D_L)$ is virtually cyclic, and all the isometries in $\Aut(\D_L)$ of infinite order have positive entropy \cite[Theorem~5.5.8]{ratcliffe}. Symmetry groups of hyperbolic lattices of rank $\le 2$ are either finite or elementary of hyperbolic type (see \cite{galluzzi.lombardo.peters}). More precisely, a hyperbolic lattice of rank $2$ has finite symmetry group if and only if it does not represent $0$ or $-2$ \cite[Corollary~3.4]{galluzzi.lombardo.peters}. We collect in the following lemma the known properties of hyperbolic lattices with an elementary symmetry group of hyperbolic type.

\begin{lemma}\label{lem:detect_virtually_abelian}
Let $L$ be a hyperbolic lattice with an elementary symmetry group of hyperbolic type. Then:
\begin{enumerate}
\item[(1)] The rank of $L$ is at most $4$.
\item[(2)] If $0\ne v\in L$ is an isotropic vector, its stabilizer in $\Aut(\D_L)$ is finite and $v^\perp/\langle v \rangle$ is a root lattice. 
%\item[(3)] There exists a negative definite primitive sublattice $N \subseteq  L$ of co-rank $2$ which is preserved by $\Aut(\D_L)$.
\item[(3)] A finite index sublattice of $L$ is spanned by $(-2)$-roots $\{\delta_i\}_{i=1,\ldots,\rk(L)}$ such that $-2 \le \delta_i . \delta_j \le 18$ for all $i,j$.
\end{enumerate}
\end{lemma}
\begin{proof}\hfill
\begin{enumerate}
\item[(1)] It follows by \cite[Theorem~9.1.1]{nikulin.factor.groups} and the preceding discussion (also cf. \cite[Theorem~3.9]{brandhorst.mezzedimi}).
\item[(2)] It follows by \cite[Theorem~3.9~and~Proposition~3.2]{brandhorst.mezzedimi} that $v^\perp/ \langle v \rangle$ is an overlattice of a root lattice. Since it has rank at most $2$, and the root lattices of rank $2$ $A_1,A_2,A_1\oplus A_1$ have no non-trivial overlattices, it is a root lattice.
% \item[(3)] Since $\Aut(\D_L)$ is elementary of hyperbolic type, it preserves a unique geodesic $\gamma$ in hyperbolic space, connecting the two fixed points in $\partial \HH_L$.
% Let $M$ be the real span of $\gamma$ in $L\otimes \RR$. Since $\gamma$ is in the positive cone, $M$ is of signature $(1,1)$. Let $f\in \Aut(\D_L)$ be of infinite order. Its characteristic polynomial is of the form $S(x)\cdot C(x)\in \ZZ[x]$, with $C(x)$ a product of cyclotomic polynomials and $S$ a Salem polynomial \Gi{Why is this? Why can't the characteristic polynomial be a Salem polynomial of degree $4$?}\Si{aposteriori it is ... But a priori, I do not see a reason right now.} Observe that $M=\ker S(f)$ is defined over the integers, so $N:=M^\perp$ is the desired lattice.
\item[(3)] This is \cite[Lemma~1]{nikulin.interesting}. \qedhere
\end{enumerate}
\end{proof}

% \begin{remark}
% Let $L$ be a hyperbolic lattice of rank $3$ or $4$ with $\Aut(L)$ elementary of hyperbolic type, $N\subseteq L$ the negative definite lattice in \Cref{lem:detect_virtually_abelian} and $M=N^\perp$.
% We have an exact sequence
% $$ \Or(N) \to \Aut(\mathcal{D}_L) \stackrel{r}{\to} \Aut(\mathcal{D}_M).$$
% \Si{I think that $\ker r$ is s subgroup of $\Or^+(N)$ of finite index, but they are not necessarily equal.} \Gi{I agree, that is why there is no $0$ on the left of the exact sequence. Do you think it is misleading written like this?}\Si{I think it is wrong as written. Or am I missing something? $O(N) \to \Aut(\D_L)$ should be given by $g \mapsto \id_M \oplus g$ ... but this does not necessarily extend to an isometry of $L$ and therefore does not define an element of its subgroup $\Aut(\D_L)$ (or factor group depending on how you want to see it). }
% By \cite[Corollary~3.4]{galluzzi.lombardo.peters}, $\Aut(\mathcal{D}_M)$ is isomorphic to either $\ZZ$ or $D_\infty = \ZZ \rtimes \ZZ/2\ZZ$, so the image of the restriction homomorphism $r$ is a subgroup of finite index of $D_\infty$. Moreover, the homomorphism $r$ splits, since any symmetry of $M$ in the image of $r$ can be extended to $L$ by acting as the identity on $N$ \Si{why can we extend by the identity on $N$? There should be some condition coming from preserving $L/(M \oplus N)$.} \Gi{You are right. I think that the real reason is that the image of $r$ is isomorphic to either $\ZZ$ or $D_\infty = \ZZ/2 * \ZZ/2$, and any surjective homomorphism to these groups splits}. On the other hand, the isometry group of a definite lattice of rank $\le 2$ embeds into the dihedral groups $D_4$ or $D_6$ (with the first group achieved by the square lattice, i.e. $A_1\oplus A_1$, and the second one by the hexagonal lattice, i.e. $A_2$). Any isometry in the kernel of $r$ commutes with the symmetries of $L$ coming from symmetries of $M$ \Si{with the overlattice buisness this is not true anymore.}, so $\Aut(\mathcal{D}_L)$ is a finite index subgroup of
% $$D_\infty \times D_4 \quad \text{or} \quad D_\infty\times D_6.$$ \Si{perhaps the proucts should be semi-direct?} \Gi{I think they are direct products, because they each act either on $M$ or $N$, so they commute}\Si{I am only convinced for $L=M \oplus N$.} \Gi{But commutativity can e.g. be checked after tensoring with $\QQ$, so I don't think there are overlattice issues.}\Si{Please give a clean proof and I will check it :-).}
% \end{remark}


We conclude this section with a crucial proposition, which is a refinement of \cite[Theorem~3.6]{mezzedimi.entropy}. We say that two groups $G_1,G_2$ are \emph{commensurable} if they admit isomorphic finite index subgroups $H_1\subseteq G_1$, $H_2\subseteq G_2$.

\begin{proposition}\label{prop:easy_fundamental_chamber}
Let $L$ be a hyperbolic lattice and $M\subseteq L$ a primitive sublattice of corank $1$. If $\disc(L)\ge 2\disc(M)$, then $M_\RR \cap \D_{L}$ is a Weyl chamber of $M$. In particular $\Aut(\D_L)$ has a subgroup commensurable with $\Aut(\D_M)$.
\end{proposition}
\begin{proof}
The fundamental chamber $\D_L$ of $L$ is cut out by hyperplanes given by the $(-2)$-roots of $L$, and therefore $M_\RR \cap \D_{L}$ is cut out by the hyperplanes $H_r \coloneqq r^\perp\otimes \RR$ given by $(-2)$-roots $r\in L$ of the form $r=x+y$, with $x\in M^\vee$ and $y\in N^\vee$, where $N\coloneqq M^\perp$. Note that $N$ is negative definite.
Clearly $M_\RR \cap \D_L$ is contained in a Weyl chamber of $M$, which we denote by $\D_M$.
If $x^2 \ge 0$, then $H_r$ does not intersect $\D_{M}^\circ$, since $M$ is a hyperbolic lattice. Therefore we may assume $x^2 < 0$, or equivalently $y^2>-2$.

In turn, the Weyl chamber of $M$ is cut out by a subset of such hyperplanes, namely those corresponding to $(-2)$-roots of $M$, or equivalently those with $y=0$. In particular $M_\RR \cap \D_{L}$ coincides with the Weyl chamber $\D_M$ if there is no $(-2)$-root $r=x+y\in L$ with $x\in M^\vee$, $y\in N^\vee$, and $y^2 \in (-2,0)$.

By contradiction, assume that there is such a $(-2)$-root $r=x+y$. The vector $y\in N^\vee$ belongs to the image $\pi(L)$ of the orthogonal projection $\pi: L \to N^\vee$. Set $n\coloneqq [L:M\oplus N] = [\pi(L):N]$, and let $-2k$ be the norm of the generator of the rank $1$ lattice $N$. We have \[\disc(L)\cdot n^2 = \disc(M)\disc(N)=2k \disc(M),\] so
$$-2 < y^2 \le \max(\pi(L)) = -2k/n^2  = -\frac{\disc(L)}{\disc(M)},$$
which implies that $\disc(L)< 2\disc(M)$, a contradiction.

For the second statement, observe that there exists a finite index subgroup $G$ of $\Or(M)$ of isometries that extend to $L$. Then, by the first part, the injective map $G\hookrightarrow \Or(L)$ descends to an inclusion $G \cap \Aut(\D_M)\hookrightarrow \Aut(\D_L)$, giving the desired claim.
\end{proof}

The following lemma will allow us to easily apply \Cref{prop:easy_fundamental_chamber} in explicit examples.

\begin{lemma}\label{lem:easy_vectors}
 Let $L$ be an even lattice and $v \in L$ primitive. Set $k \ZZ = v.L$, $kg=|v^2|$ and $M=v^\perp \subseteq L$.
 Then 
 \[2\disc(M) \leq \disc(L) \quad \Leftrightarrow \quad 2g \leq k.\]
 Moreover, $k \mid \disc(L)$. 
\end{lemma}
\begin{proof}
We note that $g=[L:M \oplus \ZZ v]$. Then
\[\disc(L) g^2 = |v^2|\disc(M)=kg \disc(M)\]
so 
\[2\disc(M) = 2 \disc(L)g/k \leq \disc(L)\]
if and only if
\[2 g \leq k.\]
Since $g$ divides $\disc(M)$, we get that $\disc(L)=k \disc(M)/g$, i.e. $k \mid \disc L$. 
\end{proof}



\section{The tests} \label{sec:tests}

In this section we present the tests that we use to decide whether a given hyperbolic lattice of rank $3$ or $4$ has an elementary symmetry group of hyperbolic type. The first one is the \emph{overlattice test}, and it is given by the following easy observation.

\begin{lemma}\label{lem:overlattice}
If $L$ has a virtually abelian symmetry group and $L'$ is an overlattice of $L$, then $L'$ has virtually abelian symmetry group as well.
Moreover, if $L$ has elementary symmetry group of hyperbolic type, then the symmetry group of $L'$ is either finite or elementary of hyperbolic type as well.
\end{lemma}
\begin{proof}
The subgroup $S$ of symmetries of $L'$ that preserve $L$ is of finite index in $\Aut(\mathcal{D}_{L'})$.
Since the fundamental chamber of $L'$ is contained in the fundamental chamber of $L$, we have $S \subseteq \Aut(\mathcal{D}_L)$. Therefore $S$ is virtually abelian and hence so is $\Aut(\mathcal{D}_{L'})$.
If $\Aut(\mathcal{D}_L)$ is elementary of hyperbolic type, then $S$ is either finite or elementary of hyperbolic type as well.
\end{proof}

The second test is called the \emph{sublattice test}. It is based on the following proposition, which in turn relies on \Cref{prop:easy_fundamental_chamber}.

\begin{proposition}\label{prop:rank2sublattice}
Let $L$ be a hyperbolic lattice of rank $3$, and assume that there exist two \emph{non-isometric} primitive sublattices $M_1,M_2$ of $L$ of rank $2$, not representing $0$ or $-2$, such that $\disc(L)\ge 2\disc(M_i)$ for $i=1,2$. Then the action of $\Aut(\mathcal{D}_L)$ on $\mathbb{H}_L$ is not elementary, and in particular $\Aut(\mathcal{D}_L)$ is not virtually abelian.
\end{proposition}
\begin{proof}
By assumption, the symmetry group of $M_1$ and $M_2$ is elementary of hyperbolic type, and in particular it is virtually cyclic. By \Cref{prop:easy_fundamental_chamber}, and up to replacing $M_2$ by $w(M_2)$ for some element $w \in W(L)$, we have $\D_{M_i} = M_i \otimes \RR \cap \D_{L}$.

Let $f_1,f_2$ be symmetries of infinite order of $M_1$ and $M_2$, respectively. \Cref{prop:easy_fundamental_chamber} implies that (some iterate of) $f_1$ and $f_2$ extend to symmetries of $L$, which, by abuse of notation, we still denote $f_1$ and $f_2$; we claim that the subgroup $\Gamma \coloneqq \langle f_1,f_2\rangle$ of $\Aut(\mathcal{D}_L)$ is not elementary, so that $\Aut(\mathcal{D}_L)$ is not elementary as well.

If $\Gamma$ is elementary, there exists a vector $v\in \partial \overline{\mathbb{H}}_L$ whose orbits under $f_1$ and $f_2$ are finite. The only finite orbit of $f_1$ (resp. $f_2$) is given by the two points in $\partial \overline{\mathbb{H}}_{M_1}$ (resp. $\partial \overline{\mathbb{H}}_{M_2}$). These two vectors generate $(M_1)_\RR$ (resp. $(M_2)_\RR$), hence if $\Gamma$ is elementary, then $\partial \overline{\mathbb{H}}_{M_1} = \partial \overline{\mathbb{H}}_{M_2}$ and in particular $(M_1)_\RR = (M_2)_\RR$. Since $M_1$ and $M_2$ are both primitive sublattices of $L$, this implies that $M_1 = M_2$, which is in contradiction to the assumption that they are not isometric.
\end{proof}

Let $L$ be a lattice and $M,N \subseteq L$ primitive sublattices such that $\rk L = \rk M+ \rk N$. 
Then we call 
\[M\oplus N \subseteq L\]
a primitive extension. Primitive extensions are in bijection with anti isometries
\[A_M \supseteq H_M \xrightarrow[]{\gamma} H_N \subseteq A_N.\]
where $|H_M|=|H_N|=[L:M \oplus N]$.
Alternatively, one may construct them from the chain of finite index inclusions 
\[M\oplus N \subseteq L \subseteq L^\vee \subseteq M^\vee \oplus N^\vee.\]
\begingroup
\captionof{algorithm}{The sublattice test}\label{alg:sublattice}
\endgroup
\begin{algorithmic}[1]
\REQUIRE A hyperbolic lattice $L$ of rank $n$. An oracle which decides if a lattice of rank $n-1$ has finite or elementary symmetry group of hyperbolic type. 
\ENSURE If the output is ``false'', then $\Aut(\mathcal{D}_L)$ is not virtually abelian.
 \STATE $\mathcal{M} = []$
 \FOR{$k \mid \disc(L)$}
  \FOR{$g \in \{g : 2g \leq k, 2 \mid gk \}$}
   \STATE Let $E$ be the set of primitive extensions
   \[[-gk]\oplus L(-1) \subseteq G\] 
   of index $k$.
   \FOR{$G \in E$}
     \STATE $D \gets A_G(-1)$
     \STATE Check if there exists a lattice $M$ of signature $(1,n-2)$ with discriminant group $A_M \cong D$. If not, continue with the next $G$, otherwise compute one such $M$.
     \STATE Compute a primitive extension $M \oplus G \subseteq \Lambda$ with $\Lambda$ unimodular.
     \STATE $L' \gets L(-1)^\perp \subseteq \Lambda$
     \STATE Check if $L \cong L'$. If not, continue with the next $G$.  
     \STATE Using the oracle, check if $\Aut(M)$ is finite or elementary of hyperbolic type. 
     \IF{$\Aut(\mathcal{D}_M)$ is not finite or hyperbolic}
       \RETURN false \label{line:not finite or hyperbolic}
     \ENDIF
     \IF{$n=3$ and $M$ does not represent $0,-2$}\label{line:M-represents-0-2}
       \IF{$M$ is not isometric to a lattice in $\mathcal{M}$}
         \STATE append $M$ to $\mathcal{M}$.
     \ENDIF
       \IF{$|\mathcal{M}|=2$ }
         \RETURN false \label{line:two-test}
       \ENDIF
    \ENDIF
    \ENDFOR
  \ENDFOR
\ENDFOR
\RETURN true
\end{algorithmic}\hrulefill\\
\begin{remark}\label{rem.explain_sublattice_test}
The procedure in \Cref{alg:sublattice} computes primitive  hyperbolic sublattices $M$ of corank $1$ of $L$ with \[2\disc(M)\le \disc(L).\] Indeed, let $M$ be such a sublattice. Let $M^\perp = \langle v \rangle$, and set $k\ZZ = v.L$, $kg = |v^2|$. By \Cref{lem:easy_vectors} it follows that $k\mid \disc(L)$ and $2g \le k$, and $L$ is a primitive extension of $M\oplus [ -gk]$ of index $g$. The lattices $L$ and $L(-1)$ glue to a unimodular lattice, hence there is a primitive extension 
\[L \oplus L(-1) \subseteq \Lambda \cong U^{n}.\] 
Hence, under the composition \[ M \oplus [-gk] \oplus L(-1) \hookrightarrow L \oplus L(-1)\hookrightarrow \Lambda,\] $M$ is the orthogonal complement of a primitive extension $G$ of $[-gk]\oplus L(-1)$ of index $k$. 
\end{remark}

\begin{lemma}
\Cref{alg:sublattice} is correct.
\end{lemma}
\begin{proof}
The iteration over $k \mid \disc(L)$ and $g$ is justified by \Cref{lem:easy_vectors} (although strictly speaking not relevant for the correctness of the algorithm).
By construction (cf. \Cref{rem.explain_sublattice_test}), we have that $M$ embeds primitively in $L$, with orthogonal complement $M^\perp \cong [-gk]$. By  
\Cref{lem:easy_vectors},
\[2 \disc(M) \leq \disc(L).\]
By \Cref{prop:easy_fundamental_chamber} the fundamental chamber $\D_M$ of $M$ is contained in some fundamental chamber $\D_L$ of $L$. Then $\Aut(M)\times \{\id_{M^\perp}\}$ is commensurable with $\Aut(L)$. Thus if $\Aut(M)$ is not finite or elementary of hyperbolic type, then neither is $\Aut(L)$.
Thus the return value in line \ref{line:not finite or hyperbolic} is correct. 

After passing line \ref{line:M-represents-0-2}, $M$ does not represent $0,-2$ and therefore $\Aut(\mathcal{D}_M)$ is elementary of hyperbolic type. Then line \ref{line:two-test} is justified by \Cref{prop:rank2sublattice}.
\end{proof}



\begin{remark}
For $n=2$ the symmetry group is always finite or elementary of hyperbolic type. 
The oracle for $n>3$ can be obtained by induction as a finite list of lattices. 
The computation of $E$ is the most expensive step. It suffices to compute it up to the action of $\Or([-gk])\times \Or(L^\vee/L(-1))$. This may or may not be faster than enumerating the primitive extensions straight away. 
\end{remark}

The third test, called the \emph{cusp test}, checks whether the lattice $L$ admits primitive isotropic vectors with infinite stabilizer. If this is the case, then the symmetry group of $L$ is not elementary of hyperbolic type by \Cref{lem:detect_virtually_abelian}. Recall that a primitive isotropic vector $v\in L$ has infinite stabilizer if and only if the lattice $M=v^\perp/\langle v \rangle$ is not a root overlattice, i.e. $\rk(M_{root}) < \rk(M)$ (cf. \cite[Proposition~3.2]{brandhorst.mezzedimi}). 

\begingroup
\captionof{algorithm}{The cusp test}\label{alg:cusp test}
\endgroup
\begin{algorithmic}[1]
\REQUIRE $L$ a hyperbolic lattice of rank $n$.
\ENSURE If the output is "false", then $\Aut(L)$ is neither finite nor elementary of hyperbolic type.
\IF{$L$ is anisotropic}
\RETURN true
\ENDIF
\STATE compute a (random) isotropic vector $v \in L$ and set $M=v^\perp/[v]$
\STATE Let $R$ be the root sublattice of $M$.
\IF{$\rk R < n -2$}
\RETURN false 
\ENDIF 
\RETURN true
\end{algorithmic}\hrulefill\\

\section{The classification} \label{sec:classification}
We are now ready to proceed with the classification of hyperbolic lattices with elementary symmetry group of hyperbolic type.
We denote the isometry class of a lattice $L$ by $[L]$.
\begingroup
\captionof{algorithm}{Classification algorithm}\label{alg:classification}
\endgroup
\begin{algorithmic}[1]
\REQUIRE An integer $n=3,4$ and an oracle to determine if $\Aut(M)$ is finite or elementary of hyperbolic type for any hyperbolic lattice $M$ of rank $n-1$.
\ENSURE A complete list of representatives of the isomorphism classes of lattices of rank $n$ with finite or elementary symmetry group of hyperbolic type.
\IF{$n=3$}  
\STATE Let $L_{a,b,c}$ denote the lattice with Gram matrix 
\[G_{a,b,c}=\begin{pmatrix}
    -2 &a &b\\
    a  &-2 &c\\
    b  &c & -2
\end{pmatrix}\]
$H_1 \gets \{ [L_{a,b,c}] : -1 \leq a \leq 2, -1 \leq b \leq c \leq 18, L_{a,b,c} \text{ is hyperbolic} \}.$ \label{alg: def H1 rk 3}
\ENDIF
\IF{$n=4$}
\STATE Let $L_{a,b,c,d,e,f}$ denote the lattice with Gram matrix 
\[G_{a,b,c,d,e,f}=\begin{pmatrix}
    -2 & a &  b &  c\\
    a  &-2 &  d &  e\\
    b  & d & -2 &  f\\
    c  & e &  f & -2\\
\end{pmatrix}\]
\STATE $H_1 \gets \{ [L_{a,b,c,d,e,f}] \text{ hyperbolic}: \  -1 \leq a \leq 2, \ a \leq b,c,d,e,f \leq 18, \ b\leq c \leq d, \ b\leq e \}$ \label{alg: def H1 rk 4}
\ENDIF
\STATE $M_1 \gets \{[M] \mid L \subseteq M \text{ is a maximal even overlattice, } [L] \in H_1\}$ \label{alg:def M1}
\STATE $M_2 \gets \{[M] \in M_1 : M \text{ passes the cusp and sublattice tests} \}$ \label{alg:def M2}
\STATE $M_3 \gets [\;]$ 
\FOR{$[M] \in M_2$}
  \STATE Compute $\Aut(M)$ using Borcherds' method 
  \IF{$\Aut(M)$ is finite or elementary of hyperbolic type }
    \STATE append $M$ to $M_3$ 
  \ENDIF
\ENDFOR
\STATE $H_2 \gets \{[L_1] : L \subseteq L_1 \subseteq M$ where $[L] \in H_1$ and $[M] \in M_3\}$ \label{alg:def H2}
\STATE sort $H_2$ by discriminant in ascending order and remove all maximal lattices from it. \label{alg:sort H2}
\STATE $H_3 \gets M_3$ \label{alg:def H3}
\FOR{$L \in H_2$} \label{alg:loop H2}
\FOR{$p \mid \disc(L)$ with $p^2 \mid \disc(L)$} \label{alg:loop indexp}
\STATE compute an even overlattice $L \subseteq L'$ of index $p$ if it exists, otherwise continue 
the loop in line \ref{alg:loop indexp} with the next prime $p$ 
\IF{$[L'] \notin H_3$}\label{alg: overlattice test}
\STATE continue the loop in line \ref{alg:loop H2} with the next lattice $L$. 
\ENDIF 
\ENDFOR
\IF{$L$ does not pass the cusp or sublattice tests}
\STATE continue the loop in line \ref{alg:loop H2} with the next lattice $L$ 
\ENDIF
\STATE Compute $\Aut(L)$ using Borcherds' method 
  \IF{$\Aut(L)$ is finite or elementary of hyperbolic type }
    \STATE append $L$ to $H_3$ 
  \ENDIF
\ENDFOR
\end{algorithmic}\hrulefill\\
\begin{lemma}
\Cref{alg:classification} is correct.
\end{lemma}
\begin{proof}
Let $L$ be a lattice with finite or elementary symmetry group of hyperbolic type. 
If $n = \rk L =3$, by \cite[Lemma 1]{nikulin.interesting} there are $3$ roots $\delta_1,\delta_2,\delta_3 \in L$ which are linearly independent and such that \[-2 < \delta_i.\delta_j \leq 18.\] 
By inspection of the proof of this lemma, we can moreover assume that $a = \delta_1.\delta_2 \in \{-1,0,1,2\}$. Denote the Gram matrix of $(\delta_1,\delta_2,\delta_3)$ by $G_{a,b,c}$
with $-1 \leq a\leq 2$ and $-1\leq b,c \leq 18$. After possibly swapping $\delta_1$ and $\delta_2$, we may assume that $b \leq c$. Let $L_{a,b,c}$ denote the lattice with Gram matrix $G_{a,b,c}$. 

If $n= \rk L =4$, by \cite[Lemma 1]{nikulin.interesting} there are $4$ roots and by inspection of the proof we can once more assume that $-1 \leq \delta_i.\delta_j \leq 2$ for at least one pair $i,j$ 
%\Si{In case 1 we take $\delta_1,\delta_2$ such that the corresponding hyperplanes meet in the point $A_2$, in case 2 we take as basis $\delta,e_1,e_2,e_3$ where the $e_i$ generate a rank $3$ root lattice.} \Gi{I agree}. 
After permuting $\delta_1,\dots,\delta_4$, we may assume that $a = \min(a,b,c,d,e,f)$. 
Then we may still exchange $\delta_1$ and $\delta_2$ as well as $\delta_3,\delta_4$ and achieve that 
$b \leq c,d,e$ as well as $c \leq d$.
By construction, the set $H_1$ contains a sublattice of $L$ of finite index.

In line \ref{alg:def H2}, $M_3$ is, by construction, the set of isometry classes of maximal even lattices with finite or elementary symmetry group of hyperbolic type. 
Therefore, $[L]$ is contained in $H_2$.

By \Cref{lem:overlattice} every overlattice $L'$ of $L$ has finite or elementary symmetry group of hyperbolic type. 
By induction and since $H_2$ is sorted by discriminant, the isometry class $[L']$ is already contained in $H_3$. Therefore, $L$ passes the overlattice test in line \ref{alg: overlattice test}.

The rest of the algorithm is clear. 
\end{proof}

\begin{remark}
A crucial part of this algorithm is to test lattices for isometry in order to create sets of isometry classes. See the appendix for an algorithm to do this. Testing the $n(n-1)/2$-pairs of a list of length $n$ for isometry is obviously too expensive if $n$ is in the range of millions. Therefore, we implemented a hash function for an isometry class that captures some invariants. The genus symbol \cite[Chapter 15]{conway.sloane.lattices} is a good invariant for this purpose. At the odd primes it is canonical and at the even ones, we simply take the ranks, scales and parity of the Jordan constituents, which is an invariant as well.  
\end{remark}

We obtain the following results from running \Cref{alg:classification}. We denote by $U(n,k)$ the rank $2$ hyperbolic lattice with Gram matrix
$$\begin{pmatrix}
    0 &n\\
    n &2k
\end{pmatrix}$$
and $U(n)\coloneqq U(n,0)$.

\begin{theorem} \label{thm:classification.3}
    Let $L$ be an even hyperbolic lattice of rank $3$ whose symmetry group is elementary of hyperbolic type. Then $L$ is isometric to exactly one of the following $45$ lattices
    \[[2n] \oplus A_2 \qquad n \in \{5, 8, 10, 11, 14, 15, 20, 24, 26, 30, 42, 48, 60, 90\}\]
    
    \[[2n]\oplus A_1 \oplus A_1 \qquad n \in \{6, 7, 12, 15, 24, 36\}\]
    
    \[\begin{pmatrix}
        2l & 1 & 0\\
        1 & -2 & 1\\
        0 & 1 & -2\\
    \end{pmatrix} \qquad l\in \{11, 13, 15, 21, 23, 33\}\]
    
     \[\begin{pmatrix}
        2l & 1 & 0\\
        1 & -2 & 0\\
        0 & 0 & -2\\
    \end{pmatrix} \qquad l\in \{5,12\}\] % 8 isometric to one listed before
     
     \[\begin{pmatrix}
        2l & 1 & 1\\
        1 & -2 & 0\\
        1 & 0 & -2\\
    \end{pmatrix} \qquad l\in \{3,7\}\] % 11,19 is isometric to another one listed before
    \[U(n) \oplus [-2] \qquad n \in \{11,14,15,20,24\}\]

    \[U(n,k)\oplus [-2] \qquad (n,k) \in \{(8,2),(8,3),(8,6),(9,3),(12,3),(12,6)(12,8),(16,8)\}\]

    \[\begin{pmatrix}
        -2&3&3 \\
        3& -2 &2 \\
        3&  2& -2 
    \end{pmatrix},\begin{pmatrix}
        -2&3&7 \\
        3& -2 &2 \\
        7&  2& -2 
    \end{pmatrix}\]
\end{theorem}
\begin{remark}
Let $n=3$. The set $H_1$ consists of $403$ isometry classes.
The sets $M_1,M_2,M_3$ consist of $192,24,21$ isometry classes.
The set $H_2$ consists of $225$ isometry classes. 
Of these, $59$ are discarded by the overlattice test in Line \ref{alg: overlattice test}, $15$ by the cusp test, $62$ by the sublattice test and, finally, $18$ by Borcherds' method.
The resulting set $H_3$ has $71$ elements. Of these $26$ have finite symmetry group, confirming Nikulin's result \cite{nikulin.finite.aut.3}, and $45$ have elementary symmetry group of hyperbolic type. See \cite{intermediate.data} for the code and intermediate results.
\end{remark}

\begin{theorem} \label{thm:classification.4}
Let $L$ be a hyperbolic lattice of rank $4$ with elementary symmetry group of hyperbolic type. 
Then $L$ is isometric to exactly one of the following $3$ lattices
\[U(6) \oplus A_1 \oplus A_1, 
\begin{pmatrix} 
0 & 4 & 0 & 0\\
 4 & -6 & 1 & 1\\
 0 & 1 & -2 & 0\\
 0 & 1 & 0 & -2
\end{pmatrix},
\begin{pmatrix} 
0 & 5 & 0 & 0\\
 5 & -6 & 2 & 2\\
 0 & 2 & -2 & 1\\
 0 & 2 & 1 & -2
\end{pmatrix}
\]
of discriminants $2^4 \cdot 3^2$, $2^6$ and $3 \cdot 5^2$.
\end{theorem} 
\begin{remark}
Let $n=4$. 
The set $H_1$ consists of $62019$ isometry classes coming from $1,533,532$ Gram matrices.  
The set $M_1$ consist of $25278$ isometry classes. After discarding $23313$ with the cusp test and $1957$ with the sublattice test, $|M_2|=8$ remain. Of these, $|M_3|=4$ pass Borcherds' method. 

The set $H_2$ consists of $2568$ isometry classes,
of these $217$, $85$ and $21$ pass the overlattice, cusp and then sublattice test. After applying Borcherds' method, $|H_3|=17$ remain. Of these, $14$ have finite symmetry group, as in \cite{vinberg.finite.aut.4}. See \cite{intermediate.data} for the code and intermediate results.
\end{remark}


\begin{remark}
We obtain a list of $48$ lattices with elementary hyperbolic symmetry group in rank $3$ and $4$. Clearly, all $3$ lattices in rank $4$ are isotropic. On the other hand, observe that one of the lattices $L$ of rank $3$ in \Cref{thm:classification.3} is isotropic if and only if $2\disc(L)$ is a square. Indeed, if $v\in L$ is isotropic, then $v^\perp/\langle v \rangle \cong A_1$ by \Cref{lem:detect_virtually_abelian}, so $\disc(L) = 2n^2$, where $n=v.L$ is the divisibility of $v$. Conversely, if $r\in L$ is a $(-2)$-root, then $\disc(r^\perp)$ is a square and thus it is represents zero. A simple computation shows that $20$ out of the $45$ lattices of rank $3$ in \Cref{thm:classification.3} are isotropic.
\end{remark}


\appendix

\section{An algorithm for weak approximation and an isometry test}\label{appendix}
In this section we give an algorithm to determine if two non-degenerate indefinite integer lattices of rank at
least $3$ are isometric. \\ 

The difficulty is the following: in order to decide if two lattices $L_1$ and $L_2$ in the same genus are in the same spinor genus (and therefore isometric), one has to produce an isometry 
$f\colon L_1 \to L_2 \otimes \QQ$ such that $[L_2:f(L_1) \cap L_2]$ is coprime to $2\det(L_1)$.
\Cite{chan.gao.weak.approx} shows that finding such a proper isometry $f$ can be done by an exhaustive search.
Following \cite[101:7]{meara.quadratic}, our approach is based on splitting local isometries into reflections and approximating them globally.

This allows to effectively decide isometry of any two integer lattices:
The case of rank $1$ is trivial and rank $2$ is solved by Gauss' reduction of binary forms. The  definite case is checked with \cite{plesken.souvignier}. Finally, two degenerate lattices are isometric if and only if they have the same rank (as abelian groups) and $L/\ker(L) \cong L'/\ker(L')$.

The algorithms are implemented as part of the computer algebra system OSCAR \cite{OSCAR}.\\


Recall that, for $V$ a rational quadratic space and $x \in V$ with $x^2 \neq 0$, the reflection in $x$ is defined by 
$s_x(y) = y- \frac{2(x.y)}{x^2} x$.
For the estimates, we use the normalized $p$-adic absolute value with $|p|_p=1/p$ and the maximum norm $|| \cdot ||_p$ on endomorphisms, matrices and vectors. The normalized $p$-adic valuation is denoted by $\nu_p(x)$ and satisfies $\nu_p(p)=1$.
If confusion is unlikely, we drop $p$ from the indices.

We recall the following definition from \cite[Section 4]{brandhorst.veniani.hensel-lifting}:
\begin{definition}
Let \(f\colon L \to M\) be a linear map of \(\ZZ_p\)-lattices of the same rank. 
We say that \(f\) is \emph{compatible} if \(f(L \cap p^iL^\vee)=M \cap p^iM^\vee\) for all \(i \in \ZZ\).
For an integer \(a \geq 0\), we call a compatible linear map \(f\) \emph{\(a\)-approximate} if there exists an isometry \(\tilde f \colon L \to M\) such that \(\tilde f \equiv f \bmod p^a\mathrm{Hom}(L,M)\). 
\end{definition}
If $L$ and $M$ are \(\ZZ_p\)-lattices and 
$f\colon L \to M\otimes \QQ_p$ is a $\ZZ_p$-linear map, then we say that $f$ is $a$-approximate if $f\colon L \to f(L)$ is $a$-approximate. Conditions to be $a$-approximate in terms of matrices are spelled out in \cite[Definition 4.2]{brandhorst.veniani.hensel-lifting}.

\begingroup
\captionof{algorithm}{Weak approximation}\label{alg:weak approximation}
\endgroup
\hrule
\begin{algorithmic}[1]
\REQUIRE \hfill
\begin{itemize}
    \item A non-degenerate quadratic space $V=\QQ^n$; 
    \item a finite set of primes $P$;
\item a collection of natural numbers $(t_p)_{p \in P}$;
\item a collection of matrices $(f_p)_{p \in P}$ with $f_p\in \mathrm{GL}(V \otimes \QQ_p)$ 
such that $f_p$ is $(\nu_p(d_p)+t_p)$-approximate, where $d_p$ is the denominator of $f_p$ and $\nu_p(\det(f_p)-1)\geq 1-\nu_p(4)$.
\end{itemize}
\ENSURE Return a matrix $f \in \mathrm{SO}(V)$ with $||f - f_p||_p \leq p^{-t_p}$ for all $p \in P$.
\FOR{$p \in P$}\label{alg:weak approximation-lfirstloop}
  \STATE $s \gets t_p+\nu_p(d_p)$
  \WHILE{true} 
  \STATE Use the algorithm described in \cite[Lemma 5.8]{shimada.connected.components} to (try to) decompose $f_p$ as a series of $p$-adic reflections, i.e. compute $x_{1,p},\dots ,x_{r_p,p} \in \ZZ^n$ such that $f_p \equiv s_{1,p} \circ \dots \circ s_{r_p,p} \pmod{ p^{t_p}}$ where $s_{i,p}$ denotes the reflection in $x_{i,p}$. \label{alg:weak approximation-lshimada} If the decomposition was successful, break the while loop and continue with the for loop in Line \ref{alg:weak approximation-lfirstloop} with the next prime. 
  \STATE $s \gets 2s$
  \STATE Use quadratic Hensel lifting of \cite[Algorithm 3]{brandhorst.veniani.hensel-lifting} to compute a matrix $f_p' \in \QQ^{n \times n}$ such that $f_p'$ is $s$-approximate and $f_p \equiv f_p' \mod p^s f_p \ZZ^{n \times n}$.  
  \STATE $f_p \gets f_p'$
  \ENDWHILE
\ENDFOR
\STATE Pick some anisotropic vector $v \in V$.
\STATE For $r_p<i\leq r:=\max\{r_p : p \in P\}$ set $x_{i,p}=v$. \label{alg:weak-approximation.padding}
\STATE $(l_i \gets 2n+10 \colon$  for $i = 1,\dots n)$
\REPEAT
    \STATE For each $1 \leq i \leq r$ compute $x_i \in \ZZ^n$ with \[x_i \equiv x_{i,p} \pmod{p^{t_p+l_i}} \quad ( \forall p \in P).\]
    \vspace*{-4mm}
    \STATE $f\gets s_{x_1} \circ \dots \circ s_{x_n}$
    \FOR{ $p \in P$}
        \STATE $e_p \gets \log_p(||f - f_p||^{-1})$
        \IF{$e_p < t_p$}
        \STATE $l_i\gets l_i+10$
        \ENDIF
    \ENDFOR
\UNTIL{$e_p \geq t_p$ for all $p \in P$}
\RETURN $f$
\end{algorithmic}\hrulefill\\

\begin{proposition}
\Cref{alg:weak approximation} terminates and is correct. 
\end{proposition}
\begin{proof}
If $f_p$ is $a$-approximate and $a$ is sufficiently high, then the decomposition into reflections in Line \ref{alg:weak approximation-lshimada} is successful. 
Since $f_p$ is $a$-approximate, it is a valid input for the Hensel lifting algorithm.
From $s \geq t_p + \nu_p(d_p)$ we get $p^sf_p \ZZ^{n \times n} \subseteq p^{t_p}\ZZ^{n\times n}$. In particular, $||f_p-f_p'||_p \leq p^{-t_p}$.   

Thus, after the for loop in Line \ref{alg:weak approximation-lfirstloop}, $||f_p - s_{x_1,p} \circ \dots \circ s_{x_{r_p},p}||_p\leq p^{-t_p}$ holds for each $p \in P$.  
Since $\nu_p(\det(f_p)-1)\geq 1 + \nu_p(4)$, the number of reflections $r_p$ is even. Therefore, we only composed with the identity in Line \ref{alg:weak-approximation.padding} and 
\[||f_p - s_{x_1,p} \circ \dots \circ s_{x_{r},p}||\leq p^{-t_p} \quad (\forall p \in P)\]
holds. 

Once the $x_i$ approximate the $x_{i,p}$ sufficiently well, the $s_{x_i}$ approximate the $s_{x_i,p}$ because the dependence of $s_{x_i}$ on $x_i$ is continuous. Since matrix multiplication is continuous, the product of the $s_{x_i}$ eventually approximates $f_p$ and the algorithm terminates. 
\end{proof}

For the sake of completeness we give an apriori estimate for the precisions $l_i$ needed for the Chinese remainder. In practice, it will not be optimal. Hence, we included the while loop.

\begin{lemma}\label{lem:error}
Let $V=\QQ^n$ be a quadratic space with Gram matrix $G$ and $h,x \in V$.
Let
\[ c = ||2G|| \cdot  |x^2|^{-1} \cdot ||x|| \max \left\{ |x^2|^{-1} \cdot ||G||, 1\right\}\]
and 
\[d = \max\{1,  ||2 G x|| /|x^2|\}\]
If $||h|| < ||x||$, then 
$ || s_x - s_{x+h}|| \leq c \cdot ||h||$
and 
\(||s_x|| = ||s_{x+h}|| \leq d. \)
\end{lemma}
\begin{proof}
\[\begin{array}{ll}
    & s_{x}(y) - s_{x+h}(y) = \dfrac{2 (x+h,y)}{(x+h)^2}(x+h) - \dfrac{2 (x,y)}{x^2}x\\[15pt]
    =& 2\left(\dfrac{1}{(x+h)^2}-\dfrac{1}{x^2} \right)(x,y) x
     + 2\dfrac{(h,y)}{(x+h)^2}x    + 2\dfrac{(h,y)}{(x+h)^2}h
     + 2\dfrac{(x,y)}{(x+h)^2}h 
\end{array}\]
We may estimate the error norm by the maximum over each summand and use that $|(x,y)|\leq ||G|| \cdot ||x|| \cdot ||y||$. 
Then we arrive at 
\[||2G|| \cdot ||y|| \cdot  |(x+h)^2|^{-1} \cdot \max \left\{ |x^2|^{-1}|(x+h)^2 - x^2| , ||h|| \cdot ||x|| , ||h||^2\right\},\]
which simplifies to
\[||2G|| \cdot ||y|| \cdot  |(x+h)^2|^{-1} \cdot \max \left\{ |x^2|^{-1}|2x.h - h^2| , ||h|| \cdot ||x|| , ||h||^2\right\}\]
and is bounded above by
\[||2G|| \cdot ||y|| \cdot  |(x+h)^2|^{-1} \cdot ||h|| \max \left\{ |x^2|^{-1}||2x|| \cdot ||G||, |x^2|^{-1} ||G|| \cdot ||h|| , ||x|| , ||h||\right\}.\]
Using the assumption $||h|| < ||x||$ we arrive at the claim.
\end{proof}
\begin{lemma}
Fix a prime $p$. Let $x_1,\dots, x_n, h_1,\dots, h_n  \in \QQ^n$ with $||h_i|| < ||x_i||$ for all $i \in \{1,\dots n \}$.
Define $c_i$ and $d_i$ as in \cref{lem:error} with $x=x_i$ and $h_i$
\[|| s_{x_1+h_1} \circ \dots \circ s_{x_n+h_n} - s_{x_1}\circ \dots \circ s_{x_n}|| \leq  \max\left\{ \prod_{i \in I} c_i ||h_i|| \prod_{j \in I^c} d_j \mid \emptyset \neq I \subseteq \{1,\dots n\}\right\} \leq Ce \]
where $e = \max\{||h_i|| : i \in \{1,\dots n\}\}$ 
and $C = \max\{\max\{c_i, d_i\} : i \in \{1,\dots , n\}\}$.
\end{lemma}
\begin{proof}
We write $s_{x_i+h} = s_{x_i}+r_i$ with $r_i = s_{x_i+h}-s_{x_i}$ in the endomorphism ring and note that $||s_{x_i}|| \leq d_i$ and $||r_i|| \leq c_i ||h_i||$. Then expand the product as a sum and estimate each term using the triangle inequality.  
\end{proof}



Given a $p$-adic lattice $L$, we say that $x \in L$ is a \emph{norm generator}, if $x^2\ZZ = n(L):=\sum_{y \in L} y^2 \ZZ$, which is called the \emph{norm} of $L$. The scale of $L$ is the (fractional) ideal $s(L):=\{x.y : x ,y \in L\}$. Since 
$2s(L)\subseteq n(L)$, the reflection in a norm generator preserves the lattice. 

\begingroup
\captionof{algorithm}{Isometry test}\label{alg:isometry-test}
\endgroup
\hrule
\begin{algorithmic}[1]
\REQUIRE $L_1,L_2$ indefinite integral lattices of rank $n \geq 3$.
\ENSURE Return whether $L_1$ and $L_2$ are isometric.
\IF{$L_1$ and $L_2$ are not in the same genus}
\RETURN false 
\ENDIF
\IF{the genus of $L_1$ consists of a single spinor genus}
\RETURN true 
\ENDIF 
\STATE Compute an isometry $f \colon L_1\otimes \QQ \to L_2\otimes \QQ$
\STATE $L_1 \gets f(L_1)$
\STATE $r \gets [L_1 : L_1 \cap L_2]$
\STATE $P = \text{support}(2 \det(L_1))$
\STATE Let $V:= L_2 \otimes \QQ = \QQ^n$, denote by $B_1$ the basis matrix of $L_1 \subseteq V$.
\STATE Let $d$ be the denominator of $B_1$.
\STATE $k_p \gets  \nu_p(r) + 2n + 10$
\FOR{$p \in P$} \label{alg:isometry-test-line-redo}
\STATE Compute a matrix $F_p \in \QQ$ such that $F_p$ is $k_p+\nu_p(d)$-approximate (w.r.t. $V$) and $F_p B_1 \in \ZZ_{(p)}$. This can be done using a $p$-adic normal form as described in \cite{miranda.morrison.embeddings} and quadratic Hensel lifting \cite[Algorithm 3]{brandhorst.veniani.hensel-lifting}. 
\IF{$\nu_p(\det(F_pB_1)-1) < 1 + \nu_p(4)$}
\STATE Compose $F_p$ by a reflection in a norm generator of $L_1 \otimes \ZZ_p$.
\ENDIF 
\ENDFOR
\STATE Compute an isometry $F \in \mathrm{SO}(V)$ with $F\equiv F_p \mod p^{\nu_p(r)}$ for all $p \in P$ using the weak approximation \Cref{alg:weak approximation}. 
\STATE $r \gets [L_2 : F(L_1) \cap L_2]$
\RETURN Return whether $r$ is an improper spinor generator of $L_1$ (see e.g. \cite[Chapter 15 \S 9]{conway.sloane.lattices}).
\end{algorithmic}\hrulefill\\

\begin{remark}
The correctness of \Cref{alg:isometry-test} follows from the theory of genera and spinor genera, see e.g. \cite{conway.sloane.lattices,cassels.quadratic,meara.quadratic}. The computation of the genus and the improper spinor generators is fast and effective, provided the factorization of the determinant is known. It is available for instance in \cite{OSCAR}.
\end{remark}


\printbibliography

\end{document}